% 06-safety.tex - Section 6: Safety Interface

\chapter{Safety Interface}
\label{sec:safety}

\section{Overview}

The NDP subsystem includes several safety features to protect equipment from damage due to temperature anomalies or hardware failures.

\section{Temperature Monitoring}

The NDP continuously monitors temperature across all subsystems:
\begin{itemize}
    \item F-engine board (ZCU102) FPGA die temperatures
    \item X-engine server temperatures (via IPMI sensors)
    \item Headnode temperatures
\end{itemize}

\subsection{Temperature Sensors}

\subsubsection{F-engine Boards}

Each ZCU102 board reports FPGA die temperature through the Xilinx system monitor interface.

\subsubsection{Head Node and X-engine Servers}

Server temperatures are monitored via IPMI sensors, which may include:
\begin{itemize}
    \item CPU temperatures
    \item GPU temperatures
    \item Ambient temperatures
    \item Power supply temperatures
\end{itemize}

\subsection{Temperature MIB Entries}

The following MIB entries report temperature status.  Table~\ref{tab:temp-mib} lists the available temperature MIB entries.

\begin{table}[htbp]
    \centering
    \caption{Temperature Monitoring MIB Entries}
    \label{tab:temp-mib}
    \begin{tabular}{lp{8cm}}
        \toprule
        \textbf{MIB Entry} & \textbf{Description} \\
        \midrule
        \mib{BOARD\{n\}\_TEMP\_MAX} & Maximum temperature on F-engine board \{n\} \\
        \mib{BOARD\{n\}\_TEMP\_MIN} & Minimum temperature on F-engine board \{n\} \\
        \mib{BOARD\{n\}\_TEMP\_AVG} & Average temperature on F-engine board \{n\} \\
        \mib{SERVER\{n\}\_TEMP\_MAX} & Maximum temperature on server \{n\} \\
        \mib{SERVER\{n\}\_TEMP\_MIN} & Minimum temperature on server \{n\} \\
        \mib{SERVER\{n\}\_TEMP\_AVG} & Average temperature on server \{n\} \\
        \mib{GLOBAL\_TEMP\_MAX} & Maximum temperature across all systems \\
        \mib{GLOBAL\_TEMP\_MIN} & Minimum temperature across all systems \\
        \mib{GLOBAL\_TEMP\_AVG} & Average temperature across all systems \\
        \bottomrule
    \end{tabular}
\end{table}

\section{IPMI Power Control}

The X-engine servers support out-of-band power control via IPMI.  Table~\ref{tab:ipmi-power} lists the supported power operations.

\begin{table}[htbp]
    \centering
    \caption{IPMI Power Operations}
    \label{tab:ipmi-power}
    \begin{tabular}{ll}
        \toprule
        \textbf{Operation} & \textbf{Description} \\
        \midrule
        status & Query current power state (on/off) \\
        on & Power on server \\
        off & Power off server \\
        cycle & Power cycle server \\
        reset & Reset server \\
        \bottomrule
    \end{tabular}
\end{table}

\section{Subsystem State Machine}

The NDP subsystem operates in one of the states shown in Table~\ref{tab:state-machine}.

\begin{table}[htbp]
    \centering
    \caption{NDP State Machine}
    \label{tab:state-machine}
    \begin{tabular}{lp{10cm}}
        \toprule
        \textbf{State} & \textbf{Description} \\
        \midrule
        \texttt{SHUTDWN} & System is powered down or not initialized. Only \cmd{INI} commands are processed. \\
        \texttt{BOOTING} & System is initializing (programming FPGAs, starting pipelines). Commands are rejected with status \hex{0C}. \\
        \texttt{NORMAL} & System is operating normally. All commands are accepted. \\
        \texttt{WARNING} & A warning condition exists but the system continues to operate. \\
        \texttt{ERROR} & An error condition exists. The system may need to be re-initialized. \\
        \bottomrule
    \end{tabular}
\end{table}

\subsection{State Transitions}

\begin{center}
\begin{tabular}{lcl}
    \texttt{SHUTDWN} & $\xrightarrow{\text{INI}}$ & \texttt{BOOTING} \\
    \texttt{BOOTING} & $\xrightarrow{\text{success}}$ & \texttt{NORMAL} \\
    \texttt{BOOTING} & $\xrightarrow{\text{failure}}$ & \texttt{ERROR} \\
    \texttt{NORMAL} & $\xrightarrow{\text{SHT}}$ & \texttt{SHUTDWN} \\
    \texttt{NORMAL} & $\xrightarrow{\text{warning}}$ & \texttt{WARNING} \\
    \texttt{NORMAL} & $\xrightarrow{\text{error}}$ & \texttt{ERROR} \\
    \texttt{WARNING} & $\xrightarrow{\text{resolved}}$ & \texttt{NORMAL} \\
    \texttt{WARNING} & $\xrightarrow{\text{error}}$ & \texttt{ERROR} \\
    \texttt{ERROR} & $\xrightarrow{\text{SHT}}$ & \texttt{SHUTDWN} \\
\end{tabular}
\end{center}

\section{F-engine Board Monitoring}

The NDP monitors F-engine board health including:
\begin{itemize}
    \item FPGA programming status
    \item FFT overflow counters
    \item Equalizer clip counts
    \item Ethernet transmit overflow/full status
    \item Data transmission status
    \item Input RMS levels
\end{itemize}

If an F-engine board fails to program or stops transmitting data, the system enters the \texttt{ERROR} state.

\subsection{Input RMS Monitoring}

In addition to the per-input statistics reported through the \mib{ANT\{n\}\_RMS} MIB entries, the NDP evaluates the distribution of input RMS values on each F-engine board to detect ADC problems.  A board is considered reasonable when both of the following hold across the inputs on that board:

\begin{itemize}
    \item The maximum RMS is less than 250 counts
    \item The standard deviation of the RMS values is less than 100 counts
\end{itemize}

The two criteria are intended to distinguish an analog signal path problem, where the gain is wrong across all inputs and the RMS values remain similar to one another, from an NDP problem, where only the inputs on one or more misbehaving ADCs depart from the rest.

\begin{calloutBox}{Operational Note}{orange}
This check is a heuristic and can mis-attribute the source of a problem.  An analog signal path fault that affects only part of the array can trip the standard deviation criterion and be reported as an NDP condition, and conversely an ADC fault common to all of the inputs on a board can present as a uniform gain error.  The reported condition should be treated as an indication that the input levels are wrong, not as a definitive statement of which subsystem is at fault.  Confirm the source against the per-input \mib{ANT\{n\}\_RMS} values and the current ASP configuration before acting.
\end{calloutBox}

The check is evaluated once per monitoring poll.  To avoid reacting to transient conditions, the system enters the \texttt{ERROR} state only after the check fails on three consecutive polls, and the condition is reported in the \texttt{SUMMARY} field of the MIB with exit code \hex{0D}:

\begin{verbatim}
    Found N FPGA board(s) with unreasonable input RMS values
\end{verbatim}

\noindent The hostnames of the affected boards are written to the system log.

\section{X-engine Pipeline Monitoring}

The NDP monitors X-engine pipeline status via ZMQ communication:
\begin{itemize}
    \item Pipeline running status
    \item Data capture status
    \item Processing errors
\end{itemize}

\section{Timing Monitor Monitoring}

The NDP monitors the timing monitor board including:
\begin{itemize}
    \item Power supply voltages (12V, 9V, 6V)
    \item Valon synthesizer lock status
    \item Recent detection of a sync pulse
    \item Microcontroller temperature
\end{itemize}

If the timing monitor board reports a problem, the system enters the \texttt{ERROR} state.

\section{Error Recovery}

To recover from most error conditions:
\begin{enumerate}
    \item Identify the underlying cause via \mib{INFO} and \mib{LASTLOG} MIB entries
    \item Address the hardware or configuration issue
    \item Issue a \cmd{SHT} command if the system is not already shut down
    \item Issue an \cmd{INI} command to re-initialize
\end{enumerate}

\begin{calloutBox}{Warning}{orange}
If a blocking operation (\cmd{INI} or \cmd{SHT}) is in progress, additional \cmd{INI} or \cmd{SHT} commands will be rejected with exit code \hex{0C} (blocking operation in progress).
\end{calloutBox}

\section{Command Validation}

The NDP validates all commands before execution:
\begin{itemize}
    \item Beam numbers must be in range 1--4
    \item Tuning numbers must be in range 1--2
    \item DRX filter codes must be in range 1--7; TBS filter codes must be in range 7--9
    \item Gain values must be in range 0--15
    \item Subslot values must be in range 0--99
\end{itemize}

Invalid commands are rejected with appropriate exit codes (see Appendix~\ref{app:codes}).
